% Deutsche Parallelfassung zu dynamics/part_intro.tex

\chapter*{Dynamik}
\phantomsection
\addcontentsline{toc}{chapter}{Dynamik}

Das identifizierte Gerade--Übergangsbogen--Kreisbogen-Alignment endet nicht an
der Grenze von Teil V. An jeder intrinsischen Position \(s\) bestimmen sein
Krümmungsgesetz, die realisierte Pose, das Profil, das Überhöhungsgesetz und
das qualifizierte bewegte Frame bereits den hier benötigten Eisenbahnzustand.
Teil VI führt deshalb kein zweites Zustandsmodell ein. Er stellt eine neue
Frage:

\begin{quote}
Wie erzeugt die Durchfahrt durch diesen räumlichen Zustand zeitabhängige
Größen, und welche zusätzlichen Modelle werden benötigt, bevor eine
Engineering-Folge bewertet werden kann?
\end{quote}

Die neue Eingabe ist ein Durchfahrtsgesetz
\[
s=s(t),\qquad v(t)=\frac{\dd s}{\dd t}.
\]
Komposition macht aus räumlichen Gesetzen zeitliche Exposition:
\(\kappa(s)\mapsto\kappa(s(t))\) und \(u(s)\mapsto u(s(t))\).
Keines der Gesetze ändert seine Identität. Es ändert sich die Rate, mit der ein
Fahrzeug oder anderes bewegtes System sie durchfährt.

\begin{figure}[htbp]
\centering
\input{figures/dynamics/space_to_time_dependency-DE}
\caption{Vom qualifizierten Eisenbahnzustand über zeitliche Exposition und Antwortmodell zur Bewertung und möglichen Entscheidung.}
\label{fig:space-to-time-dependency}
\end{figure}

Jeder Pfeil in \cref{fig:space-to-time-dependency} ergänzt Information:
Das Durchfahrtsgesetz ergänzt Zeit, seine Ableitungen Geschwindigkeit und
Beschleunigung, eine reduzierte Beziehung angegebene kinematische Annahmen,
ein Antwortmodell Fahrzeug- und Gleisphysik und die Bewertung Zweck und
anwendbare Regeln. Erst eine Engineering Decision ergänzt eine autorisierte
Festlegung.

\begin{table}[htbp]
\centering
\small
\begin{tabular}{p{0.19\textwidth}p{0.33\textwidth}p{0.37\textwidth}}
\toprule
\textbf{Ebene} & \textbf{Erforderliche Eingabe} & \textbf{Ergebnis und Grenze}\\
\midrule
Räumliches Gesetz & identifiziertes Alignment, \(s,\kappa(s),u(s)\) & intrinsische Geometrie und getrenntes Überhöhungsgesetz; keine Zeit\\
Durchfahrtsgesetz & \(s(t)\), Quelle, Anwendbarkeit, Gültigkeit & \(v(t)\), \(\dot v(t)\), zeitliche Exposition; keine Fahrzeugantwort\\
Reduzierte Kinematik & Frame, Vorzeichen, Überhöhungs- und Geschwindigkeitsannahmen & Indikatoren wie \(v^2\kappa\); nicht Komfort oder Sicherheit\\
Antwortmodell & Fahrzeug-, Kontakt-, Gleis- und Lastdaten & modellabhängige Antwort; weder Beobachtung noch Entscheidung\\
Bewertung & Zweck, Regel, Geltungsbereich, Evidenz & qualifizierte Beurteilung; ohne eigene Autorität\\
Entscheidung & zuständige Autorität und zugelassene Evidenz & angenommener oder abgelehnter Handlungsweg\\
\bottomrule
\end{tabular}
\caption{Der Dynamikargumentgang ergänzt jeweils genau eine Verantwortung.}
\label{tab:dynamics-responsibility-layers}
\end{table}

\begin{principle}[Prinzip der dynamischen Zustandsinterpretation]
\label{princ:dynamic-state-interpretation}
Zeitabhängige Größen werden durch Komposition des qualifizierten
Eisenbahnzustands mit einem qualifizierten Durchfahrtsgesetz und, soweit
erforderlich, expliziten Antwortmodellen abgeleitet. Reduzierte Kinematik,
Antwort, Bewertung und Entscheidung bleiben getrennt.
\end{principle}

\begin{summary}
Geometrie liefert das räumliche Gesetz. Durchfahrt ergänzt Zeit. Modelle
liefern Antwort. Zweck und Regeln liefern Bewertung. Autorität liefert
Entscheidung.
\end{summary}
